\documentclass[12pt]{article}
\usepackage[utf8]{inputenc}
\usepackage[T1]{fontenc}
\usepackage[margin=1in]{geometry}
\usepackage{setspace}
\usepackage{xcolor}
\usepackage[colorlinks=true, urlcolor=blue!60!black, linkcolor=blue!60!black]{hyperref}
\setlength{\parindent}{0em}
\setlength{\parskip}{0.45em}
\renewcommand{\baselinestretch}{1.0}
\begin{document}
\begin{center}
{\Large \textit{Arizona v. United States}, 567 U.S. 387 (2012)\par}
\vspace{0.3em}
{\normalsize Case Brief --- One-Page Edition\par}
\end{center}

\textbf{Memory Jogger.} Most of Arizona's S.B.~1070 is preempted by federal immigration law, but the ``show me your papers'' check survives a facial challenge (\href{https://supreme.justia.com/cases/federal/us/567/387/}{Justia}).

\textbf{Facts.} Pursuing ``attrition through enforcement,'' Arizona made it a state crime to fail to carry registration papers (\S~3) or for an unauthorized alien to seek work (\S~5(C)), authorized warrantless arrests for removable offenses (\S~6), and required status checks during lawful stops (\S~2(B)). The United States sued on preemption grounds.

\textbf{Procedural History.} The district court enjoined all four provisions; a divided Ninth Circuit affirmed.

\textbf{Votes.} 5--3 (Kagan recused): Kennedy for the Court; Scalia, Thomas, and Alito each concurred in part and dissented in part.

\textbf{Holding.} Sections 3, 5(C), and 6 are preempted by field or obstacle preemption; Section 2(B) is not facially preempted before state courts construe it, if it does not prolong detention solely to check status.

\textbf{Reasoning.} Congress occupies the field of alien registration (\textit{Hines}); IRCA deliberately imposed only civil consequences on unauthorized workers; Section 6 gave state police more arrest power than federal officers have. Congress encourages information sharing, so Section 2(B) could await construction. Scalia invoked the States' inherent power to exclude; Thomas would let only enacted text preempt; Alito would have upheld Sections 5(C) and 6.

\textbf{Examples.} \emph{Same-Side:} a state crime of unlawful entry. \emph{Opposite-Side:} a neutral W-4 identity-theft prosecution (\textit{Kansas v. Garcia}). \emph{Fence-Sitter:} mandatory 48-hour ICE detainer holds.

\textbf{Critique.} Critics fault the Court both for ignoring racial profiling and for inferring preemption from congressional silence and executive discretion.

\textbf{Key Quotation.} ``The Government of the United States has broad, undoubted power over the subject of immigration and the status of aliens.'' --- Justice Kennedy, 567 U.S. at 394.

\end{document}
